%
\documentclass[11pt]{article}
\usepackage{xspace,mathrsfs}
\usepackage{graphicx,rotating}
\usepackage{url}
\usepackage{amscd,amsfonts,amsmath,amssymb}
\usepackage{color}
\usepackage{longtable}
\usepackage{multirow,setspace}
\renewcommand{\topfraction}{0.9}
\renewcommand{\textfraction}{0.1}
\usepackage{float}
\usepackage{booktabs}
\usepackage{enumerate}
\usepackage{caption}
\usepackage{fancyhdr}

\def\shownotes{0}   % set 1 for version with author notes
                    % set 0 for no notes

%%%%%%%%%%%%%%%%%%%%%%%%%%%%%%%%%%%%%%%%%%%%%%%%%%%%%%%%%%%%%%
%Some macros (you can ignore everything until "end of macros")
\topmargin 0pt \advance \topmargin by -\headheight \advance
\topmargin by -\headsep

\textheight 8.9in
\oddsidemargin 0pt \evensidemargin \oddsidemargin \marginparwidth
0.5in

\textwidth 6.5in

%%%%%%
\DeclareMathAlphabet{\mathsl}{OT1}{cmr}{m}{sl}
\DeclareMathAlphabet{\mathsc}{OT1}{cmr}{m}{sc}
\DeclareMathAlphabet{\mathslbf}{OT1}{cmr}{bx}{sl}

\newcommand{\procfont}[1]{\mathsc{#1}}
\newcommand{\tablefont}[1]{\mathsf{#1}}
\newcommand{\Initialize}{\procfont{Initialize}}
\newcommand{\Finalize}{\procfont{Finalize}}
\newcommand{\FinalizeCodeNoarg}{\underline{\textrm{proc} $\Finalize$}}
\newcommand{\Onquery}[2]{\underline{\textrm{proc} $#1(#2)$}}
\newcommand{\InitializeCode}{\underline{\textrm{proc} $\Initialize$}}
\newcommand{\FinalizeCode}[1]{\underline{\textrm{proc} $\Finalize(#1)$}}

\newcommand{\fprff}{\algfont{f}}
\newcommand{\prim}{\primfont{P}}
\newcommand{\adv}{\advfont{Adv}}

\newcommand{\gamefont}[1]{\mathrm{#1}}

\newcommand{\getsr}{{{\,\leftarrow{\hspace*{-3pt}\raisebox{.75pt}{$\scriptscriptstyle\$$}}\,}}}
\newcommand{\bits}{\{0,1\}}
\newcommand{\bit}{\{0,1\}}
\newcommand{\Ex}{\mathbb{E}}
\newcommand{\eqdef}{\stackrel{def}{=}}
\newcommand{\To}{\rightarrow}
\newcommand{\e}{\epsilon}
\newcommand{\R}{\mathbb{R}}
\newcommand{\N}{\mathbb{N}}
\newcommand{\Z}{\mathbb{Z}}
\newcommand{\ie}{\emph{i.e.,} }
\newcommand{\eg}{\emph{e.g.,} }
\newcommand{\VAR}{\mathrm{VAR}}
\newcommand{\Prob}[1]{{\Pr\left[\,{#1}\,\right]}}
\newcommand{\probb}[2]{{\Pr}_{#1}\left[\,{#2}\,\right]}
\newcommand{\probbs}[3]{{\Pr}^{#1}_{#2}\left[\,{#3}\,\right]}
\newcommand{\probbp}[2]{{\Pr}_{#1}'\left[\,{#2}\,\right]}
\newcommand{\Probb}[2]{{{\Pr}_{#1}\left[\,{#2}\,\right]}}
\newcommand{\Probbp}[2]{{{\Pr}_{#1}'\left[\,{#2}\,\right]}}
\newcommand{\condProb}[2]{{\Pr}\left[\,#1\,|\,#2\,\right]}
\newcommand{\expect}[1]{{{\mathbf{E}}[{#1}]}}
\newcommand{\next}{ \: ; \: }
\newcommand{\true}{\mathsf{true}}
\newcommand{\false}{\mathsf{false}}
\newcommand{\AND}{\;\wedge\;}
\newcommand{\bad}{\mathsf{bad}}
\newcommand{\MsgSp}{\mathsf{MsgSp}}
\newcommand{\Gm}{\textnormal{G}}
\newcommand{\fn}{\footnotesize}
\newcommand{\schemefont}[1]{{\mathsf{#1}}}
\newcommand{\algfont}[1]{{\mathsf{#1}}}
\newcommand{\Adv}{\mathbf{Adv}}
\newcommand{\AdvGen}[3]{\mathbf{Adv}^{\mathrm{#1}}_{#2}(#3)}
\newcommand{\Exp}{\mathbf{Exp}}

\newcommand{\INDCPA}{\mbox{IND-CPA}}
\newcommand{\AdvINDCPA}[2]{\AdvGen{ind\text{-}cpa}{#1}{#2}}
\newcommand{\AdvINTCTXT}[2]{\AdvGen{int\text{-}ctxt}{#1}{#2}}

\newcommand{\LR}{\mathsf{LR}}

\newcommand{\calA}{{\cal A}}
\newcommand{\calB}{{\cal B}}
\newcommand{\calC}{{\cal C}}
\newcommand{\calF}{{\cal F}}
\newcommand{\calK}{{\cal K}}
\newcommand{\calM}{{\cal M}}
\newcommand{\calO}{{\cal O}}
\newcommand{\calR}{{\cal R}}
\newcommand{\calH}{{\cal H}}
\newcommand{\calG}{{\cal G}}
\newcommand{\calD}{{\cal D}}
\newcommand{\calE}{{\cal E}}
\newcommand{\calP}{{\cal P}}
\newcommand{\calS}{{\cal S}}
\newcommand{\calT}{{\cal T}}
\newcommand{\calU}{{\cal U}}
\newcommand{\calX}{{\cal X}}
\newcommand{\calV}{{\cal V}}
\newcommand{\calW}{{\cal W}}
\newcommand{\calY}{{\cal Y}}
\newcommand{\calN}{{\cal N}}

\DeclareMathAlphabet{\mathsl}{OT1}{cmr}{m}{sl}

\providecommand{\mypara}[1]{\smallskip\noindent\emph{#1} }
\providecommand{\myparab}[1]{\smallskip\noindent\textbf{#1} }
\providecommand{\myparasc}[1]{\smallskip\noindent\textsc{#1} }
\providecommand{\para}{\smallskip\noindent}

\newcommand{\Keys}{\mathsf{Keys}}
\newcommand{\Dom}{\mathsf{Dom}}
\newcommand{\Rng}{\mathsf{Rng}}

\newcommand{\main}{\procfont{Game}}

\newcommand{\gamesfontsize}{\small}
\renewcommand{\captionfont}{}
\newcommand{\indd}{\hspace*{12pt}}

\newcommand{\mpage}[2]{\begin{minipage}{#1\textwidth} #2 \end{minipage}}
\newcommand{\fpage}[2]{\framebox{\begin{minipage}[t]{#1\textwidth}\setstretch{1.2}\gamesfontsize  #2 \end{minipage}}}

\newcommand{\hfpages}[3]{\hfpagess{#1}{#1}{#2}{#3}}
\newcommand{\hfpagess}[4]{
        \begin{tabular}{c@{\hspace*{.5em}}c}
        \framebox{\begin{minipage}[t]{#1\textwidth}\setstretch{1.2}\gamesfontsize #3 \end{minipage}}
        &
        \framebox{\begin{minipage}[t]{#2\textwidth}\setstretch{1.2}\gamesfontsize #4 \end{minipage}}
        \end{tabular}
    }

\newcommand{\twoCols}[4]{
\begin{center}
        \framebox{
        \begin{tabular}{c@{\hspace*{.4em}}|c@{\hspace*{.4em}}c}
        \begin{minipage}[t]{#1\textwidth}\setstretch{1.2}\gamesfontsize #3 \end{minipage}
        &
        \begin{minipage}[t]{#2\textwidth}\setstretch{1.2}\gamesfontsize #4 \end{minipage}
        \end{tabular}
        }
\end{center}
}

\newcommand{\gameName}[1]{\underline{$\main$ #1} \\}
\newcommand{\procName}[1]{\underline{#1} \\[2pt]}

\newcommand{\bea}{\begin{eqnarray*}}
\newcommand{\eea}{\end{eqnarray*}}

\newcommand{\return}{\textrm{Ret}~}

\newtheorem{thm}{Theorem}[section]
\newtheorem{lem}[thm]{Lemma}
\newtheorem{cor}[thm]{Corollary}
\newtheorem{propo}[thm]{Proposition}
\newtheorem{clm}[thm]{Claim}
\newtheorem{defn}[thm]{Definition}
\newtheorem{assm}[thm]{Assumption}
\newtheorem{rem}[thm]{Remark}
\newtheorem{obs}[thm]{Observation}
\newtheorem{egs}[thm]{Example}
\newtheorem{fct}[thm]{Fact}
\newtheorem{expr}{Experiment}
\newtheorem{cons}[thm]{Construction}
\newtheorem{nte}[thm]{Note}
\newtheorem{coj}[thm]{Conjecture}
\newtheorem{ques}[thm]{Question}

%%%%%%%  Author Notes %%%%%%%
\ifnum\shownotes=1
\newcommand{\authnote}[2]{{ $\ll$\textsf{\footnotesize Solution: #2}$\gg$}}
\else
\newcommand{\authnote}[2]{}
\fi
\newcommand{\Anote}[1]{{\authnote{Adam}{#1}}}
%%%%%%%%%%%%%%%%%%%%%%%%%%%%%%%%%

% Scheme-specific macros for this exercise
\newcommand{\GF}{\mathrm{GF}(2^{128})}
\newcommand{\GHASH}{\mathrm{GHASH}}
\newcommand{\Poly}{\mathrm{Poly1305}}
\newcommand{\ChaCha}{\mathrm{ChaCha20}}
\newcommand{\AESGCM}{\mathsf{AES\text{-}GCM}}
\newcommand{\CCP}{\mathsf{ChaCha20\text{-}Poly1305}}
\newcommand{\ctr}{\mathrm{ctr}}

% end of macros
%%%%%%%%%%%%%%%%%%%%%%%%%%%%%%%%%%%%%%%%%%%%%%%%%%%%%%%%%%%%%%


\title{CS~466: Group Exercise}

\date{}

% page counting, header/footer
\usepackage{fancyhdr}
\usepackage{lastpage}
\pagestyle{fancy}
\lhead{\footnotesize \parbox{11cm}{CS 466, UMass Amherst, Spring 2026} }
\rhead{\footnotesize Instructor:  Adam O'Neill \\ \url{adamo@cs.umass.edu}}
\renewcommand{\headheight}{24pt}

\begin{document}

\maketitle

%Don't forget to add this for header.
\thispagestyle{fancy}

\paragraph*{GROUP EXERCISE.}

\bigskip
\noindent
In lecture we studied authenticated encryption via \emph{generic
composition}.
Here we examine $\AESGCM$ and $\CCP$---the two AEAD schemes mandated
by TLS~1.3---to see how that theory plays out in practice, and to work
through the nonce-reuse attack that breaks both schemes when nonces are
mismanaged.
All problems set $A = \varepsilon$.

\bigskip

%%%% BACKGROUND %%%%

\noindent\textbf{Nonce-Based AEAD.}\quad
\begin{sloppypar}
An \emph{authenticated encryption with associated data (AEAD)}
scheme $\mathcal{AE} = (\calK, \calE, \calD)$ consists of an
encryption algorithm
$\calE \colon \calK \times \calN \times \calA \times \calM \to \calC$
and a decryption algorithm
$\calD \colon \calK \times \calN \times \calA \times \calC \to \calM \cup \{\bot\}$,
where $\calN$ is the \emph{nonce space} and $\calA$ is the
\emph{associated-data space}.
$\calE$ is \emph{deterministic}: the caller supplies a nonce
$N \in \calN$ instead of random coins, but must \emph{never reuse}
$N$ under the same key.
Correctness: $\calD_K(N, A, \calE_K(N, A, M)) = M$.
The scheme is \emph{secure} if it achieves both IND-CPA
(ciphertext indistinguishability) and INT-CTXT (ciphertext integrity) defined below.
\end{sloppypar}

The \emph{associated data}~$A$ is authenticated but not encrypted.
Real protocols attach unencrypted metadata to every message---a TLS
record type, a sequence number, a session ID---that must be
\emph{authentic} (to prevent replay or context-confusion attacks) but
need not be secret.
AEAD authenticates $A$ and $M$ together in one pass: $\calD$ outputs
$\bot$ if $A$ is tampered with, binding each ciphertext to its context.
TLS~1.3 sets $A$ to the 5-byte record header; we set $A=\varepsilon$
below to focus on the core structure.

\smallskip
{\hfuzz=3pt\begin{center}
\hfpagess{0.42}{0.51}{%
  \gameName{IND-CPA$^b$}
  \vspace*{-\topsep}
  \begin{tabbing}
  \hspace{0.6em}\=\hspace{1.2em}\=\kill
  \> \InitializeCode\\
  \>\> $b \getsr \{0,1\}$; $K \getsr \calK$; $\mathit{asked} \gets \emptyset$\\[3pt]
  \> \Onquery{Enc}{N,A,M_0,M_1}\\
  \>\> $\textbf{if}\ |M_0| \ne |M_1|\ \textbf{then return}\ \bot$\\
  \>\> $\textbf{if}\ N \in \mathit{asked}\ \textbf{then return}\ \bot$\\
  \>\> $\mathit{asked} \gets \mathit{asked} \cup \{N\}$\\
  \>\> $\textbf{return}\ \calE_K(N,A,M_b)$\\[3pt]
  \> \FinalizeCode{b'}\\
  \>\> $\textbf{return}\ b' = b$
  \end{tabbing}
}{%
  \gameName{INT-CTXT}
  \vspace*{-\topsep}
  \begin{tabbing}
  \hspace{0.6em}\=\hspace{1.2em}\=\kill
  \> \InitializeCode\\
  \>\> $K \getsr \calK$;\; $\mathcal{S} \gets \emptyset$;\; $\mathit{win} \gets \mathbf{false}$\\[3pt]
  \> \Onquery{Enc}{N,A,M}\\
  \>\> $C \gets \calE_K(N,A,M)$\\
  \>\> $\mathcal{S} \gets \mathcal{S}\cup\{(N,A,C)\}$\\
  \>\> $\textbf{return}\ C$\\[3pt]
  \> \Onquery{Dec}{N,A,C}\\
  \>\> $\textbf{if}\ (N,A,C)\in\mathcal{S}\ \textbf{then return}\ \bot$\\
  \>\> $M \gets \calD_K(N,A,C)$\\
  \>\> $\textbf{if}\ M \ne \bot\ \textbf{then}\ \mathit{win} \gets \mathbf{true}$\\
  \>\> $\textbf{return}\ M$\\[3pt]
  \> \FinalizeCodeNoarg\\
  \>\> $\textbf{return}\ \mathit{win}$
  \end{tabbing}
}
\end{center}}

\bigskip

\noindent\textbf{Arithmetic in $\GF$.}\quad
$\GF$ (a \emph{Galois field}) is a finite number system with
$2^{128}$ elements where you can add, subtract, multiply, and divide.
Its elements are 128-bit strings; \emph{addition} is XOR
($a + b = a \oplus b$), and \emph{multiplication} is polynomial
arithmetic modulo a fixed irreducible polynomial
(hardware-accelerated via \texttt{PCLMULQDQ}).
The key fact: \emph{a nonzero polynomial of degree~$d$ over $\GF$
has at most $d$ roots}---exactly as over~$\mathbb{R}$.
We write $H^\ell$ for the $\ell$-fold product in~$\GF$.

\bigskip

\noindent\textbf{GHASH.}\quad
Since addition in $\GF$ equals XOR, $\oplus$ and $+$ are
interchangeable below.
Given hash key $H \in \bits^{128}$ and message blocks
$X[1]\|\cdots\|X[\ell]$, GHASH evaluates a polynomial at $H$:
\[
  \GHASH_H(X[1]\|\cdots\|X[\ell])
  \;=\;
  X[1] \cdot H^{\ell}
  \;\oplus\;
  X[2] \cdot H^{\ell-1}
  \;\oplus\;
  \cdots
  \;\oplus\;
  X[\ell] \cdot H.
\]
The AES-GCM tag has the form $T = \GHASH_H(C) \oplus S$,
where $S = E_K(\ctr(N,0))$.
This differs from the deterministic MACs in lecture, where security
comes from modeling $\mathrm{MAC}_K$ as a PRF\@.
Here, $\GHASH_H$ alone is not pseudorandom: $H$ is a static key, so
evaluations at the same $H$ are algebraically related.
Security instead rests on $S$: it is a fresh pseudorandom value
unknown to the adversary, making $T$ unpredictable despite the
algebraic structure of $\GHASH_H(C)$.
Crucially, $S$ must be used \emph{at most once}---reusing a nonce
repeats $S$, which is exactly what Problem~3 exploits.

\bigskip

\noindent\textbf{$\AESGCM$.}\quad
$E \colon \bits^{128} \times \bits^{128} \to \bits^{128}$ is AES-128;
$K \getsr \bits^{128}$.
Messages are block-aligned: $M[1]\|\cdots\|M[\ell]$, each block 128~bits.
$\ctr(N, i)$ denotes the 128-bit block formed from nonce $N \in \bits^{96}$
and counter~$i$.

{\hfuzz=3pt\begin{center}
\hfpagess{0.45}{0.49}{%
  \underline{\textbf{Alg} $\calE_K(N, M)$}\\
  \vspace*{-\topsep}
  \begin{tabbing}
  \hspace{0.8em}\=\hspace{1.2em}\=\kill
  \> $H \gets E_K(0^{128})$\\
  \> $\textbf{for}\ i = 1\ \textbf{to}\ \ell\ \textbf{do}$\\
  \>\> $C[i] \gets M[i] \oplus E_K(\ctr(N,i))$\\
  \> $T \gets \GHASH_H(C[1]\|\cdots\|C[\ell])$\\
  \> \phantom{$T \gets {}$}$\oplus\; E_K(\ctr(N,0))$\\
  \> $\textbf{return}\ C[1]\|\cdots\|C[\ell]\|T$
  \end{tabbing}
}{%
  \underline{\textbf{Alg} $\calD_K(N,\; C\|T)$}\\
  \vspace*{-\topsep}
  \begin{tabbing}
  \hspace{0.8em}\=\hspace{1.2em}\=\kill
  \> $H \gets E_K(0^{128})$\\
  \> $T' \gets \GHASH_H(C[1]\|\cdots\|C[\ell])$\\
  \> \phantom{$T' \gets {}$}$\oplus\; E_K(\ctr(N,0))$\\
  \> $\textbf{if}\ T \ne T'\ \textbf{then return}\ \bot$\\
  \> $\textbf{for}\ i = 1\ \textbf{to}\ \ell\ \textbf{do}$\\
  \>\> $M[i] \gets C[i] \oplus E_K(\ctr(N,i))$\\
  \> $\textbf{return}\ M[1]\|\cdots\|M[\ell]$
  \end{tabbing}
}
\end{center}}

\bigskip

\noindent\textbf{Poly1305.}\quad
Poly1305 uses the same two-part tag structure as GHASH
(polynomial hash $+$ one-time mask), but over $\mathbb{Z}_p$
for prime $p = 2^{130}-5$.
A key pair $(r, s)$ plays the same roles as $(H,\, E_K(\ctr(N,0)))$ in
AES-GCM: $r$ is the \emph{evaluation point} and $s$ is the
\emph{one-time mask}.
Each message block $X[i]$ is treated as an integer $x_i \in \mathbb{Z}_p$.
The tag is
\[
  \Poly_{r,s}(X[1]\|\cdots\|X[\ell])
  \;=\;
  \Bigl(\sum_{i=1}^{\ell} x_i \cdot r^{\ell+1-i} \bmod p\Bigr)
  + s \;\bmod 2^{128}.
\]
Since $p$ is prime, $\mathbb{Z}_p$ is a field and so the at-most-$d$-roots fact above applies.
A fresh $(r, s)$ is derived for every nonce, so $(r, s)$ is never reused.

\bigskip

\noindent\textbf{$\CCP$.}\quad
ChaCha20 is a stream cipher; $K \getsr \bits^{256}$, nonce
$N \in \bits^{96}$.
Write $\mathit{ks}_K(N, j)$ for the $j$-th 128-bit keystream block
(counters $j \ge 1$).
The Poly1305 key $(r,s)$ is the first 256~bits of $\ChaCha_K(N,0)$.

{\hfuzz=3pt\begin{center}
\hfpagess{0.45}{0.49}{%
  \underline{\textbf{Alg} $\calE_K(N, M)$}\\
  \vspace*{-\topsep}
  \begin{tabbing}
  \hspace{0.8em}\=\hspace{1.2em}\=\kill
  \> $(r, s) \gets$ first $256$ bits of $\ChaCha_K(N, 0)$\\
  \> $\textbf{for}\ j = 1\ \textbf{to}\ \ell\ \textbf{do}$\\
  \>\> $C[j] \gets M[j] \oplus \mathit{ks}_K(N, j)$\\
  \> $T \gets \Poly_{r,s}(C[1]\|\cdots\|C[\ell])$\\
  \> $\textbf{return}\ C[1]\|\cdots\|C[\ell]\|T$
  \end{tabbing}
}{%
  \underline{\textbf{Alg} $\calD_K(N,\; C\|T)$}\\
  \vspace*{-\topsep}
  \begin{tabbing}
  \hspace{0.8em}\=\hspace{1.2em}\=\kill
  \> $(r, s) \gets$ first $256$ bits of $\ChaCha_K(N, 0)$\\
  \> $T' \gets \Poly_{r,s}(C[1]\|\cdots\|C[\ell])$\\
  \> $\textbf{if}\ T \ne T'\ \textbf{then return}\ \bot$\\
  \> $\textbf{for}\ j = 1\ \textbf{to}\ \ell\ \textbf{do}$\\
  \>\> $M[j] \gets C[j] \oplus \mathit{ks}_K(N, j)$\\
  \> $\textbf{return}\ M[1]\|\cdots\|M[\ell]$
  \end{tabbing}
}
\end{center}}

\bigskip


\paragraph*{Problem 1.} (AES-GCM.)

\noindent\textbf{Part A} (5 points).
Which generic composition paradigm does $\AESGCM$ follow? Briefly justify your answer.

\medskip

\noindent\textbf{Part B} (10 points).
Consider a variant of $\AESGCM$ that replaces $E_K(\ctr(N,0))$ with
the constant $E_K(0^{128})$, making the tag mask independent of the nonce.
Show this variant is \emph{not} INT-CTXT.

\bigskip

\paragraph*{Problem 2.} (ChaCha20-Poly1305.)

\noindent\textbf{Part A} (5 points).
Which generic composition paradigm does $\CCP$ follow? Briefly justify your answer.

\bigskip

\paragraph*{Problem 3.} (Nonce-Reuse Attack on AES-GCM, 20 points.)

\noindent
Show that $\AESGCM$ is \emph{not} INT-CTXT when nonces may be reused.

\noindent\emph{Hint}: $T_1 \oplus T_2$ yields a polynomial equation in $H$ over $\GF$. Moreover, $\GF$ is a field, so division (i.e.\ multiplication by the inverse) is defined for any nonzero element.


\end{document}
